\section{Historical Timeline of Heresies and Responses}\label{sec:timeline}

This chronology is event-based. A movement can appear more than once because refutation, synodal judgment, definition, civil enforcement, reconciliation, and later ecumenical clarification are different events. An asterisk marks a primarily civil or coercive measure included to keep the aftermath visible; it is not a magisterial act. Dates prefaced by \emph{c.} are approximate. The source and qualification for every row are recorded in the repository's timeline audit.

\begin{historytimeline}
c.~49 & Jerusalem; apostles and elders & Claim that Gentile converts require circumcision and the Mosaic law for salvation & Acts 15 rejects that necessity and gives limited communion requirements; Galatians locates justification and covenant membership in Christ. & Apostolic boundary, not a trial under later canon law; grade A.\\
c.~107 & Smyrna/Rome; Ignatius of Antioch & Teachers who made Christ's flesh and suffering an appearance & Ignatius insists that Jesus was truly born, ate, suffered, died, and rose, and binds Eucharistic and ecclesial communion to that confession. & Early episcopal refutation; target school uncertain; grade B/C.\\
c.~144 & Rome; presbyterate & Marcion's rival god, edited canon, and docetic Christ & Marcion was excluded from the Roman community; Justin, Irenaeus, and Tertullian answered him by the one Creator, one economy, and apostolic Scriptures. & No extant Roman proposition decree; a durable Marcionite church followed; grade C.\\
c.~150--190 & Rome and Lyons; Justin and Irenaeus & Simonian, Menandrian, Cerinthian, Saturnilian, Basilidean, Carpocratian, Valentinian, and related systems & Apologies and \emph{Against Heresies} set witness-specific accounts beside one Creator, one incarnate Son, public apostolic faith, moral judgment, and bodily salvation. & Patristic refutations, not one council or proof that the named schools formed a single institution; grade B--D by claim.\\
c.~200 & Asia Minor; regional bishops & Montanist or New Prophecy claims and discipline & Eusebius reports synods examining and rejecting the movement; later writers contest prophecy, ecclesial authority, and rigorism. & Acts lost; original and later Montanism must be distinguished; grade C.\\
251 & Rome; Pope Cornelius and Roman synod & Novatian's rival episcopate and rigorist refusal of reconciliation after grave post-baptismal sin & The Roman synod upheld Cornelius, excluded Novatian and his party, and defended the Church's power to reconcile the penitent. & Novatian communities endured for centuries; mixed schism and ecclesiological error; grade B.\\
268 & Antioch; regional synod & Paul of Samosata's dynamic or adoptionist Christology & The synod deposed Paul and rejected a Christology in which the divine Word was not a subsisting divine person incarnate. & Imperial enforcement removed Paul only later; fragments are mediated by opponents; grade B.\\
314 & Arles; western council & Donatist challenge to Caecilian and sacramental communion & Arles rejected the Donatist appeal and treated ordination and baptism as not nullified by a minister's unworthiness. & Rival African hierarchy continued; imperial coercion later intensified; grade B.\\
325 & Nicaea I; ecumenical council & Arius's created and mutable Son & The creed confesses the Son begotten from the Father's substance, true God from true God, consubstantial with the Father; attached anathemas reject ``there was when he was not'' and creature language. & Arius exiled by imperial order*; decades of formula conflict followed; grade A/B.\\
341--351 & Antioch and Sirmium; eastern synods & Marcellian and Photinian accounts of Word, kingdom, and Incarnation & Successive synods deposed or condemned Marcellus and Photinus and issued competing creedal formulas. & Local acts became entangled with the post-Nicene party struggle; not one ecumenical settlement; grade B.\\
359--360 & Rimini, Seleucia, Constantinople; imperial councils & Homoian formula refusing substance-language & Court-backed councils imposed ``like according to the Scriptures'' and proscribed \emph{ousia} terminology. & Coercive reception was brief; these were not ecumenical definitions received by the Catholic Church; grade A/B.\\
362 & Alexandria; synod under Athanasius & Divisions among Nicenes and early denial of the Spirit's divinity & The \emph{Tomus ad Antiochenos} recognizes orthodox uses of one and three hypostases, confesses the Spirit's inseparability, and safeguards Christ's complete humanity. & Important bridge to Constantinople I; grade A.\\
377--382 & Rome and Constantinople; Damasus and Constantinople I & Eunomians, Arians, Pneumatomachi, Sabellians, Marcellians, Photinians, and Apollinarians & Roman acts and the 381 council reaffirm Nicene faith; the creed confesses the Spirit as Lord and life-giver, and canon 1 names the rejected groups. & Conciliar reception fixed the pro-Nicene Trinitarian grammar; full 381 acts do not survive; grade A/B.\\
380 & Saragossa; regional council & Ascetic practices associated with Priscillianist circles & Canons correct secret meetings, liturgical absence, fasts, and ascetic disruptions without supplying a complete dogmatic system. & Later heresiological charges exceed this act; grade B.\\
385 & Trier; imperial court* & Priscillian and companions after episcopal proceedings & The secular court executed them for \emph{maleficium}; Martin of Tours and others opposed the capital procedure. & First famous execution of Christians for heresy-related charges; civil sentence is not a creed; evidence contested.\\
c.~390--393 & Rome and Milan; Siricius and western bishops & Jovinian's propositions on baptismal security, ascetical merit, fasting, reward, and Mary's virginal childbirth & A Roman council excluded Jovinian and associates; the Milanese bishops received the judgment while defending both marriage's goodness and virginity's distinctive excellence. & No durable Jovinian communion can be traced; Jerome's polemic is not the sentence; grade B/C.\\
411 & Carthage; imperial conference* & Catholic and Donatist episcopal claims & A state-appointed official ruled for the Catholic party after a mass conference; Augustine's theology distinguished sacramental validity from fruitful communion. & Repressive laws followed; Donatism declined unevenly and survived into later centuries; grade A/B.\\
418 & Carthage and Rome; African council and Pope Zosimus & Pelagian and Caelestian claims about original sin, infant baptism, grace, and sinlessness & African canons affirm inherited sin, baptism's remission, and grace necessary not only for knowledge but for willing and doing good; Zosimus confirmed the condemnation. & Pelagius's own positions and later Pelagian systems require distinction; grade A/B.\\
431 & Ephesus; ecumenical council & Nestorius's refusal of Cyril's Christological and Marian language & The council received Cyril's second letter, condemned Nestorius, and defended one incarnate subject truly born from Mary as \emph{Theotokos}. & Rival session and political conflict followed; the 433 reunion clarified language; grade A.\\
433 & Alexandria and Antioch; Formula of Reunion & Post-Ephesine division over unity and distinction in Christ & Cyril and John of Antioch confess one Christ, perfect God and perfect man, with union without confusion and Mary as \emph{Theotokos}. & Demonstrates reconciliation through mutually interpreted formulas; grade A.\\
448--451 & Constantinople and Chalcedon; synods and ecumenical council & Eutyches's account of Christ after the union and the violence of the 449 council & Chalcedon received Leo's Tome and defined one and the same Son in two natures, without confusion, change, division, or separation; it deposed Dioscorus on conciliar-procedural grounds. & Imperial enforcement* and lasting separation followed; Eutyches is not a synonym for every non-Chalcedonian; grade A.\\
c.~473 & Arles; regional bishops & Propositions attributed to the priest Lucidus on grace, freedom, Christ's saving will, and predestination to evil & Lucidus subscribed a profession repudiating the listed claims and confessing grace's priority without divine predestination to sin. & Retraction text may reflect an accuser's construction; no ecumenical act or durable Lucidian school; grade B/C.\\
484--519 & Rome and Constantinople; Acacian schism & Reception of Chalcedon and the imperial \emph{Henotikon} & Roman and Constantinopolitan sees broke communion; the Formula of Hormisdas restored it in 519 with an explicit Chalcedonian confession. & Primarily a communion crisis with doctrinal stakes, not a new single heresy; grade A.\\
529 & Orange II; regional council confirmed by Boniface II & Gallic propositions later called ``Semi-Pelagian'' & Canons make even the beginning of faith and desire for cleansing the work of prevenient grace, reject natural unaided saving good, and also reject predestination to evil. & Retrospective movement name; positive synthesis became authoritative in the West; grade A.\\
543 & Constantinople; Emperor Justinian and eastern bishops & Origenist propositions including pre-existence and necessary restoration & An imperial edict and subscribed anathemas rejected a defined set of propositions. & Relation of the famous fifteen anathemas to the 553 council remains disputed; grade B.\\
553 & Constantinople II; ecumenical council & Theodore of Mopsuestia, specified writings of Theodoret, and Ibas's letter---the Three Chapters & The council condemned the person and writings while receiving Ephesus and Chalcedon and insisting on the unity of the incarnate Word. & Western resistance and the northern Italian schism followed; not a newly organized ``Three Chapters heresy''; grade A.\\
c.~557--567 & Alexandria and Constantinople; rival episcopal parties & Tritheist formulations treating the three hypostases as three particular natures or deities & Chalcedonian and non-Chalcedonian theologians rejected the conclusion; Alexandrian censures and Constantinopolitan negotiations preserved the one undivided divine essence. & Fragmentary, partisan record and no separate Catholic ecumenical canon against Philoponus; grade B/C.\\
c.~600 & Alexandria and Rome; Eulogius and Gregory I & Agnoete claim that Christ's humanity was ignorant of the day and hour & Eulogius refuted the claim; Gregory's \emph{Register} 10.39 endorsed the response and distinguished knowledge in Christ's humanity from knowledge derived from human nature as source. & Papal theological correspondence, not an ecumenical proposition list; named party fades; grade B/C.\\
649 & Rome; Lateran synod under Martin I & Monoenergist and Monothelite formulas and imperial \emph{Ecthesis}/\emph{Typos} & The synod confesses two natural wills and operations in the one Christ and condemns named authors and documents. & Martin was arrested and exiled by imperial power*; the teaching prepared Constantinople III; grade A.\\
680--681 & Constantinople III; ecumenical council & One will or operation in Christ & The council defines two natural wills and operations, the human will freely following the divine, and anathematizes named teachers including Honorius. & Pope Leo II confirmed the council in 682 while characterizing Honorius's fault precisely; grade A.\\
754 & Hieria; imperial council* & Defense of sacred images & The iconoclast council prohibited images and claimed conciliar authority. & Later rejected as a false council; iconophile persecution accompanied its policy; grade A for its act, not Catholic authority.\\
787 & Nicaea II; ecumenical council & Iconoclasm & The council distinguishes worship due to God from honor shown to images and teaches that honor passes to the prototype. & Reception continued through Byzantine conflict and western translation disputes; grade A.\\
792--799 & Regensburg, Frankfurt, Rome; western synods & Spanish Adoptionism of Elipandus and Felix & Western councils reject calling Christ, as man, God's adoptive son and confess the one natural Son in both natures. & Felix recanted at Rome; pockets of dispute remained in Spain; grade A/B.\\
843 & Constantinople; Byzantine restoration & Renewed iconoclasm after Nicaea II & Empress Theodora's government restored images and the synodical feast of Orthodoxy received the council. & End of official Byzantine iconoclasm, though not every western image controversy; grade A/B.\\
848--853 & Mainz, Quierzy, and Soissons; Carolingian synods & Gottschalk's double-predestination formulations & Mainz and Quierzy condemned and imprisoned Gottschalk; competing theologians and later councils articulated differing formulas on predestination and grace. & No simple universal ``Gottschalk heresy'' verdict resolves the Carolingian dispute; coercive imprisonment* continued; grade B.\\
1050--1079 & Vercelli, Rome, and Tours; papal councils & Berengar's Eucharistic account & Repeated councils rejected denials of substantial conversion and required professions, culminating in the 1079 formula of true conversion and Christ's substantial presence. & Berengar submitted; language developed toward the later term transubstantiation; grade A/B.\\
c.~1112--1124 & Utrecht and Antwerp; canons and Norbert & Propositions and a personal cult attributed to Tanchelm & Utrecht canons requested detention and judgment; Norbert later preached Catholic ministry and Eucharistic communion at Antwerp. & Accusation and pastoral response survive, but no target text or synodal sentence; Tanchelm's killing was not an extant conciliar judgment; grade C/D.\\
1139 & Lateran II; ecumenical council & Unnamed groups rejecting Eucharist, infant baptism, priesthood, and marriage & Canon 23 condemns those who feign religion while repudiating these sacramental and ecclesial goods. & Often associated with several dissenting currents, but the canon does not supply one secure sociology; grade B.\\
1140--1141 & Sens; regional council and Innocent II & Propositions attributed to Peter Abelard & Sens condemned a list concerning Trinity, grace, Christ, and sin; Innocent ordered silence and burning of the book. & Abelard appealed, then submitted and reconciled at Cluny; attribution/context remain debated; grade B.\\
1148 & Reims; council under Eugene III & Gilbert of Poitiers's Trinitarian terminology & The council required corrections distinguishing God, divinity, persons, and essence. & Gilbert submitted; no durable Gilbertine church resulted; proposition wording is mediated; grade B.\\
1163 & Tours; papal council & Heretics in southern France, often associated with dualist currents & Canonical measures prohibited sheltering them and directed secular princes to restrain them. & Pre-inquisitorial repression; nomenclature does not yet prove one transregional Cathar organization; grade B.\\
1179 & Lateran III; ecumenical council & Cathars, Patarenes, Publicani, and related names & Canon 27 anathematizes named groups and their defenders and promises spiritual privileges to armed opponents. & Joined doctrine to coercive response; labels are regional and inconsistent; grade A/B.\\
1184 & Verona; Lucius III and Frederick I & Cathars, Patarenes, Humiliati, Poor of Lyon, Passagians, Josephines, Arnoldists & \emph{Ad abolendam} lists groups, imposes episcopal inquiry, and coordinates ecclesial excommunication with secular penalties*. & Major legal framework; group listing does not mean identical doctrine; grade A.\\
1199 & Innocent III & Protection and spread of heresy & \emph{Vergentis in senium} compares pertinacious heresy to treason and provides confiscatory consequences*. & Influential penal analogy; not a new dogmatic definition; grade A.\\
1208 & Rome; Innocent III & Waldensian reconciliation & Durand of Huesca and companions profess Catholic faith, sacraments, hierarchy, and the legitimacy of oaths and temporal property. & Shows that reconciliation and differentiated treatment accompanied repression; grade A.\\
1209--1229 & Languedoc; papal and princely forces* & Albigensian/Cathar presence and political resistance & The Albigensian Crusade used warfare against territories and protectors while preaching and councils pursued doctrinal and pastoral response. & Enormous violence and political conquest cannot be treated as a doctrinal formula; regional Cathar evidence remains contested.\\
1210 & Paris; local ecclesial and royal action* & Amalricians and David of Dinant's writings & Amalrician propositions were condemned; persons were degraded and executed or imprisoned, and David's books prohibited. & Target texts fragmentary and reports hostile; local acts, grade B/D.\\
1215 & Lateran IV; ecumenical council & Cathar dualism; Joachim's criticism of Peter Lombard's Trinitarian language & Constitution 1 confesses one Creator, Incarnation, sacraments, resurrection, and judgment; constitution 2 rejects Joachim's proposition while expressly protecting his monastery and acknowledging his submission; constitution 3 legislates against heretics*. & Doctrinal precision and coercive law must be separated; grade A.\\
1231 & Gregory IX & Heresy procedure and recurring groups & \emph{Excommunicamus} consolidates canonical penalties and authorizes papally delegated inquiry. & Institutional papal inquisition developed; procedure is not itself a new creed; grade A.\\
1252 & Innocent IV and Lombard cities* & Secular prosecution of convicted heretics & \emph{Ad extirpanda} regulates civil enforcement and permits torture under stated limits. & Coercive and morally grave aftermath, not magisterial definition; grade A.\\
1255 & Anagni; Alexander IV commission & Gerard of Borgo San Donnino's Joachite \emph{Introduction to the Eternal Gospel} & A commission censured the work's radicalized age-of-the-Spirit claims; papal suppression and Franciscan discipline followed. & Distinguish Gerard's construction from everything Joachim wrote; grade B.\\
1270/1277 & Paris; Bishop Stephen Tempier & Academic propositions concerning intellect, necessity, providence, ethics, and theology & Tempier condemned 13 and later 219 propositions within his jurisdiction. & Local university censures, not 232 universal heresies; influence on scholastic debate remains contested; grade A/B.\\
1311--1312 & Vienne; ecumenical council & Errors attributed to certain Beguards/Beguines; soul-body doctrine & \emph{Ad nostrum} condemns specified perfectionist propositions; \emph{Fidei catholicae fundamento} teaches the rational soul as per se and essential form of the body. & The council did not condemn every Beguine community; grade A.\\
1317--1323 & Avignon; John XXII & Spiritual Franciscan/Fraticelli disobedience and absolute-poverty proposition & \emph{Sancta Romana} and \emph{Gloriosam Ecclesiam} condemn factions; \emph{Cum inter nonnullos} calls pertinacious assertion that Christ and the apostles owned nothing individually or in common erroneous and heretical. & Some friars submitted; others were imprisoned or executed*; distinguish poverty doctrine from movement identity; grade A.\\
1327 & Avignon; John XXII & Marsilius of Padua and John of Jandun & \emph{Licet iuxta doctrinam} rejects five conclusions subordinating or dissolving ecclesial spiritual authority. & Authors remained under imperial protection; mixed political and doctrinal controversy; grade A.\\
1329 & Avignon; John XXII & Meister Eckhart's propositions & \emph{In agro dominico} grades 17 articles as erroneous or tainted with heresy and 11 as suspect, while noting Eckhart's conditional revocation. & Later reception must preserve individual grades and the fact that not all his work was condemned; grade A/B.\\
1377 & Rome; Gregory XI & Nineteen conclusions from John Wyclif & Papal letters send propositions for examination, calling some heretical and others erroneous and dangerous. & Wyclif retained protection; later English and conciliar proceedings expanded the lists; grade A/B.\\
1382 & London; Blackfriars council & Wycliffite Eucharistic and ecclesial propositions & Archbishop Courtenay's council grades ten conclusions heretical and fourteen erroneous. & Oxford restrictions followed; Lollard reception was diverse; grade A/B.\\
1415 & Constance; council and civil authority* & Wyclif's memory and Jan Hus & Constance condemns extracted articles, orders Wyclif's remains removed, and sentences Hus after disputed attribution and recantation demands; the secular authority executes Hus. & Proposition grades and conciliar authority differ by decree/session; Hussite war and division followed; grade A/B.\\
1418 & Constance/Rome; Martin V & Wyclifite and Hussite adherence & \emph{Inter cunctas} supplies interrogation articles and describes the earlier materials collectively as heretical, erroneous, rash, seditious, or offensive. & Definitively blocks the claim that every listed article bore one censure; grade A.\\
1433--1436 & Basel/Bohemia; council parties and estates & Utraquist communion and Hussite demands & The Compactata permit a regulated communion-under-both-kinds settlement while requiring Catholic Eucharistic doctrine and ecclesial order. & Limited settlement with Utraquists; Taborite and later reception differed; disputed conciliar setting; grade A/B.\\
1460 & Rome; Pius II & Appeal from papal judgment to a future council & \emph{Execrabilis} condemns such appeals as an abuse contrary to ecclesial order. & Principal anti-conciliarist reference; jurisdictional-doctrinal act, not a catalogue of every conciliar theory; grade A.\\
1513 & Lateran V; ecumenical council & Soul mortality and one-intellect theses & \emph{Apostolici regiminis} affirms each human soul's immortality and multiplicity and requires philosophy teaching compatible with faith. & Aimed at academic propositions, not Aristotle or philosophy as such; grade A.\\
\end{historytimeline}

\clearpage
\begin{historytimeline}
1520 & Rome; Leo X & Forty-one propositions from Martin Luther's works & \emph{Exsurge Domine} condemns the propositions collectively as heretical, scandalous, false, offensive, seductive, or contrary to Catholic truth and requires retraction. & The act does not assign an individual grade to every proposition; Luther burns it; grade A.\\
1521 & Rome/Worms; Leo X and Charles V* & Luther's refusal to retract & \emph{Decet Romanum Pontificem} declares excommunication; the Edict of Worms imposes imperial outlawry*. & Ecclesial penalty and civil enforcement remain distinct; confessional separation expands; grade A.\\
1525 & Toledo; Spanish Inquisition & Forty-eight propositions judicially attributed to Alumbrados, \emph{dexados}, or \emph{perfectos} & An edict of faith states and qualifies the allegations individually, addressing abandonment, prayer, sacraments, moral effort, and claims of impeccability. & Inquisitorial construction, not one movement's confession or a universal papal definition; later individual sentences followed; grade B.\\
1525--1527 & Z\"urich and other cities; civic-church councils* & Anabaptist rejection of infant baptism and rebaptism of adults & Local disputations reject the teaching; Catholic and Protestant governments impose exile and, in some places, death*. & Anabaptists were highly diverse; coercion by several confessions is not a doctrinal proof; grade A/B.\\
1546 & Trent, sessions IV--V; ecumenical council & Reformation disputes over revelation, Scripture, tradition, original sin, and baptism & Trent receives Scripture and apostolic traditions, identifies the canonical books and authentic Vulgate use, and defines original sin and its remission without including Mary in the decree. & First conciliar phase; canons target propositions rather than naming every communion; grade A.\\
1547 & Trent, sessions VI--VII & Justification and sacraments in general; Baptism and Confirmation & The council teaches prevenient grace, faith, hope, charity, real interior renewal, merit in grace, seven sacraments instituted by Christ, and sacramental efficacy not dependent on ministerial faith. & Canonical condemnations preserve precise opposites; Catholic and Protestant readings continued to develop; grade A.\\
1551 & Trent, sessions XIII--XIV & Eucharist, Penance, and Extreme Unction & The council defines true, real, substantial presence and transubstantiation; sacramental confession and absolution; and Anointing as a true sacrament. & Zwinglian, Reformed, Lutheran, and Catholic positions must not be collapsed; grade A.\\
1553 & Geneva; civic court* & Michael Servetus's anti-Trinitarian and anti-baptismal teaching & Reformed Geneva convicts and executes Servetus after consultation with other Swiss churches. & Not a Catholic judgment; included as confessional aftermath and warning against coercion; target writings survive, grade A.\\
1562--1563 & Trent, sessions XXI--XXV & Communion, sacrifice of the Mass, Orders, Matrimony, purgatory, saints, images, indulgences & The council completes proposition-level definitions and reform decrees, distinguishing sacramental doctrine from disciplinary reform. & Becomes the Catholic confessional answer to the Reformations; grade A.\\
1567/1579 & Rome; Pius V and Gregory XIII & Seventy-nine propositions associated with Michael Baius & \emph{Ex omnibus afflictionibus} condemns the list with several censures applied collectively; \emph{Provisionis nostrae} reaffirms the judgment, especially on nature, original justice, freedom, grace, and merit. & Baius submitted with later disputes about wording and censure distribution; not every item may be called individually heretical; grade A/B.\\
1607 & Rome; Paul V & Thomist--Molinist dispute over grace and free choice & After the \emph{de auxiliis} congregations, Paul V permits both schools to teach within stated limits and forbids mutual censure. & A decisive exclusion from the heresy census: unresolved Catholic schools are not heresies.\\
1612--1613 & Sens, Aix, and Rome; provincial councils and Paul V & Edmond Richer's restriction of papal and ministerial jurisdiction & Sens censures the book's propositions \emph{as they sound} under mixed notes; Aix follows, Paul V confirms the episcopal response, and the Index prohibits the book. & No numbered papal syllabus or one-grade censure of every sentence; Richerist influence persists; grade A/B.\\
1642 & Rome; Urban VIII & Jansenius's \emph{Augustinus} & \emph{In eminenti} prohibits the book as renewing condemned Baian and grace propositions. & Publication and reception continue, leading to proposition-level judgment; grade A.\\
1653 & Rome; Innocent X & Five propositions on commandments, grace, freedom, and Christ's death & \emph{Cum occasione} condemns four as heretical and one as false or heretical according to its stated qualification. & The ``fact'' dispute---whether Jansenius taught them in the condemned sense---follows; grade A.\\
1656 & Rome; Alexander VII & Jansenist fact/right distinction & \emph{Ad sanctam beati Petri sedem} confirms that the five propositions are in \emph{Augustinus} in the condemned sense. & Formulary controversy and ``respectful silence'' persist; grade A/B.\\
1665--1679 & Rome; Alexander VII and Innocent XI & Laxist moral propositions & Roman decrees condemn successive lists permitting grave conduct through distorted probability, intention, restitution, and sacramental reasoning. & Condemned moral errors with varied notes, not a single organized heresy; grade A.\\
1682/1690 & Paris and Rome; French assembly and Alexander VIII & Four Gallican Articles on temporal power, conciliar limits, custom, and reformability of papal judgments & \emph{Inter multiplices} declares the 1682 proceedings and articles null and void; Vatican I later defines primacy and infallibility under exact conditions. & A broad ecclesial-political tendency, not four identically graded heresies; grade A.\\
1687 & Rome; Innocent XI & Sixty-eight propositions associated with Miguel de Molinos & \emph{Caelestis Pastor} condemns Quietist propositions that empty ascetic effort, petition, resistance to temptation, and ecclesial mediation in the name of passive annihilation. & Molinos abjures and is imprisoned; broader mystical theology is not condemned; grade A/B.\\
1690 & Rome; Alexander VIII & Jansenist/rigorist and moral propositions & \emph{Sanctissimus Dominus} condemns thirty-one propositions, including distortions of grace, freedom, ecclesiology, and moral action. & Demonstrates Roman rejection of both laxism and rigorism; proposition-level act, grade A.\\
1699 & Rome; Innocent XII & Twenty-three propositions from F\'enelon's \emph{Maxims of the Saints} & \emph{Cum alias} condemns the propositions on pure love and indifference to salvation as rash, scandalous, ill-sounding, offensive, dangerous, or erroneous. & F\'enelon submits; not every article is styled heresy and not all mystical disinterested love is rejected; grade A.\\
1713 & Rome; Clement XI & 101 propositions from Pasquier Quesnel & \emph{Unigenitus Dei Filius} condemns the propositions collectively under multiple grades, concerning grace, Scripture, Church, sacraments, and authority. & Appeals and Gallican/Jansenist resistance divide France; individual-grade caution required; grade A.\\
1764 & Rome/German lands; Index and Clement XIII & Febronius's reduction of papal jurisdiction and strong national-episcopal theory & The book is prohibited and Clement XIII directs the German bishops against it. & Book-level suppression, not a proposition-by-proposition dogmatic definition; Josephinism is treated as a related state policy, grade A/B.\\
1794 & Rome; Pius VI & Eighty-five propositions of the Synod of Pistoia & \emph{Auctorem fidei} assigns proposition-specific or grouped censures on Church, grace, sacraments, liturgy, and discipline. & Reform program suppressed; this act is a model for exact censure parsing, not a blanket label; grade A.\\
1832 & Rome; Gregory XVI & Religious indifferentism and an unbounded liberty-of-conscience proposition & \emph{Mirari vos} rejects indifference among religions and liberty detached from truth and the common good. & Vatican II later teaches civil immunity from religious coercion while preserving the duty to seek truth; genres and objects differ, grade A.\\
1835 & Rome; Gregory XVI & Georg Hermes's method and works & \emph{Dum acerbissimas} prohibits the works for positive doubt and rationalist treatment of faith and associated doctrines. & Mixed lower censures and book prohibition; no blanket personal-heretic verdict, grade A/B.\\
1835--1840 & Rome; Gregory XVI and Augustin Bautain & Fideist propositions on reason and revelation & Bautain signs theses affirming natural rational knowledge of God, motives of credibility, and reason's genuine capacity while requiring revelation for mysteries. & Personal submission through required propositions; philosophical fideism, not faith itself, is corrected; grade A.\\
1855 & Rome; Holy Office and Augustin Bonnetty & Radical traditionalist propositions & Required theses affirm reason's capacity to know God and demonstrate credibility while maintaining revelation's supernatural necessity. & Corrective subscription, not an act declaring every traditionalist a heretic; grade A.\\
1857 & Rome; Pius IX & Anton G\"unther's philosophical theology & \emph{Eximiam tuam} confirms the prohibition of works whose rational reconstruction endangered Trinity, creation, soul, and Incarnation. & G\"unther submits; the act does not condemn philosophical inquiry as such; grade A/B.\\
1861 & Rome; Holy Office & Seven Ontologist propositions & The decree rejects immediate knowledge of God as the necessary ground of all human cognition and related accounts of being and creation. & A precise philosophical censure; grade A.\\
1864 & Rome; Pius IX & Naturalism, religious indifferentism, coercive secular absolutism, and related errors & \emph{Quanta cura} directly condemns specified propositions; the attached \emph{Syllabus} indexes earlier papal acts across eighty headings. & Genre and original-source context control interpretation; not every modern liberty is condemned under one undifferentiated note; grade A.\\
1870 & Vatican I; ecumenical council & Pantheism, materialism, rationalism, fideism, denial of revelation's credibility, and denial of papal primacy/infallibility & \emph{Dei Filius} defines God, revelation, faith, and faith-reason harmony; \emph{Pastor aeternus} defines Roman primacy and ex cathedra infallibility under exact conditions. & Old Catholic resistance and separation follow; definitions are narrower than later polemical slogans; grade A.\\
1873--1875 & Rome; Pius IX & Old Catholic rejection of Vatican I and construction of a rival hierarchy & \emph{Etsi multa} identifies the doctrinal denials and invalid episcopal election; \emph{Graves ac diuturnae} addresses the separated and schismatic body. & Apply the council's definitions to propositions and \emph{Unitatis redintegratio} 3 to later descendants; grade A.\\
1887 & Rome; Holy Office & Forty propositions extracted from Antonio Rosmini & \emph{Post obitum} condemns the propositions as not consonant with Catholic truth. & The CDF's 2001 note later distinguishes heterodox interpretations from Rosmini's authentic thought and removes reasons for prior doctrinal concern; grade A with later qualification.\\
1899 & Rome; Leo XIII & Tendencies reported under the name ``Americanism'' & \emph{Testem benevolentiae} conditionally rejects minimizing external guidance, vows, traditional virtues, and inherited methods, while expressing confidence that American bishops did not hold the views. & Not evidence of one organized American heresy; grade A.\\
1907 & Rome; Holy Office and Pius X & Modernist propositions on revelation, dogma, Scripture, Christ, sacraments, and Church & \emph{Lamentabili} condemns sixty-five propositions; \emph{Pascendi} constructs and refutes an interconnected Modernist system; \emph{Praestantia Scripturae} confirms the acts and attaches penalties. & Individual propositions and authors still require exact attribution; ``Modernism'' is not a synonym for modern scholarship; grade A.\\
1910 & Rome; Pius X & Persistence of Modernist method among clergy and teachers & \emph{Sacrorum antistitum} imposes the anti-Modernist oath and disciplinary vigilance. & Oath later discontinued; disciplinary response must be distinguished from each doctrinal censure; grade A.\\
1937 & Rome; Pius XI & Atheistic communism as a materialist and collectivist worldview & \emph{Divini Redemptoris} rejects denial of God, soul, afterlife, personal dignity, and the class-struggle system. & A condemned rival ideology, often held by the unbaptized, not ordinarily canonical heresy under canon 751; grade A.\\
1944 & Rome; Holy Office & Mitigated or spiritual millenarianism & A decree states that a system in which Christ would visibly reign on earth before the final judgment cannot safely be taught, even in mitigated form. & Censure is ``cannot safely be taught,'' not a newly defined universal heresy name; grade A.\\
1949 & Rome; Holy Office & Support for Communist organizations and profession of materialist anti-Christian doctrine & A 1 July decree distinguishes enrollment/support/publication conduct and the profession or defense of materialist Communist doctrine. & Penal and doctrinal provisions; Communism itself is not simply a baptized Christian sect; grade A.\\
1949 & Rome/Boston; Holy Office & Restrictive reading of ``outside the Church no salvation'' associated with Leonard Feeney's circle & \emph{Suprema haec sacra} teaches the Church's necessity while affirming implicit desire ordered by supernatural faith and charity; later disciplinary action concerns refusal of obedience. & The doctrine is neither indifferentism nor a claim that only registered Catholics can be saved; grade A.\\
1950 & Rome; Pius XII & Errors concerning dogmatic development, relativism, reason, Scripture, original sin, and evolution & \emph{Humani generis} warns against specified false or dangerous opinions, permits research on human bodily origins under conditions, and excludes polygenism as then irreconcilable with received original-sin doctrine. & The encyclical does not declare evolution as such a heresy or condemn all renewed theology; grade A.\\
1953 & Rome; Holy Office & Leonard Feeney's continued refusal of obedience & A decree declares excommunication after repeated summons and contumacy. & The penalty is for grave disobedience, not a decree formally adjudging personal heresy; grade A.\\
1956 & Rome; Holy Office & Situation ethics & An instruction rejects making personal intuition or circumstances, rather than objective moral law, finally determinative. & Prohibited moral system, not a formal dogmatic definition of a named heresy; grade A.\\
1965 & Vatican II; ecumenical council & Enduring Christian divisions and historical polemic & \emph{Unitatis redintegratio} teaches that descendants are not guilty of the original separation and requires truthful representation, hierarchy of truths, and humble dialogue without false irenicism. & Positive doctrinal and pastoral qualification of how historical labels are used; grade A.\\
1966 & Rome; CDF & Postconciliar errors involving Scripture, dogma, Christ, Eucharist, penance, original sin, morality, and ecumenism & A circular letter to episcopal conferences identifies propositions requiring vigilance and correction. & Inventory of errors, not one new named heresy; grade A.\\
1967 & Rome; CDF & Profession of faith and anti-Modernist oath discipline & A new profession replaces the earlier Tridentine formula and anti-Modernist oath for specified uses. & Disciplinary replacement does not reverse the substantive doctrinal acts of 1907; grade A.\\
1973 & Rome; Paul VI and Shenouda III & Catholic--Coptic Christological inheritance & Common declaration professes one incarnate Lord, perfect God and perfect man, without mixing, commingling, confusion, alteration, division, or separation. & Historic Eutychian propositions remain rejected; present Coptic Orthodox faith is not simply labeled Eutychian; grade A.\\
1973 & Rome; CDF & Relativizing dogmatic formulas and Church indefectibility & \emph{Mysterium Ecclesiae} distinguishes historical expression and legitimate development from treating dogmatic meaning as indefinitely replaceable. & Modern doctrinal clarification, not a new named-heretic register; grade A.\\
1979 & Rome; CDF & Hans K\"ung's positions on ecclesial infallibility, Magisterium, and Eucharistic ministry & \emph{Christi Ecclesia} judges departure from integral Catholic truth and withdraws authorization to teach as a Catholic theologian. & Text/office judgment, not a formal personal-heresy trial; grade A.\\
1984/1986 & Rome; CDF & Reductive uses of Marxist analysis and Christian liberation & The two instructions reject class-struggle hermeneutics, violence, materialist anthropology, and reduction of salvation while affirming integral liberation and the preferential concern for the poor. & They do not condemn every liberation theology or struggle for justice; grade A.\\
1984 & Rome; John Paul II and Ignatius Zakka I & Catholic--Syriac Christological inheritance & Common declaration confesses the same incarnate Word and says the confusions and schisms over the Incarnation have no real basis in the faith now jointly confessed. & Requires qualified present terminology while preserving historic acts; grade A.\\
1988 & Rome; John Paul II; later status in 2009 & Lefebvre episcopal consecrations and SSPX status & \emph{Ecclesia Dei} calls the consecrations a schismatic act; Benedict XVI's 2009 \emph{Ecclesiae unitatem} states that remission of censures did not grant canonical status and doctrinal questions remained. & Canonical/communion boundary, not a declaration that the SSPX is one heresy; grade A.\\
1994 & Rome; John Paul II and Mar Dinkha IV & Catholic--Assyrian Christological inheritance & Common declaration professes one Lord Jesus Christ, true God and true man, and explains that earlier anathemas concerned persons and formulas amid misunderstanding. & ``Nestorian Church'' is rejected as an adequate current label for the Assyrian Church of the East; grade A.\\
1994/1995 & Rome; John Paul II and CDF & Reservation of priestly ordination to men & \emph{Ordinatio sacerdotalis} states the Church lacks authority and the teaching is definitively to be held; the CDF responsum identifies ordinary-universal infallibility. & Definitive doctrine; the 1998 commentary places it in the second assent category rather than automatically calling every denial technical heresy; grade A.\\
1998 & Rome; CDF & Confusion among revealed dogma, definitive doctrine, and other authentic teaching & Doctrinal commentary on the \emph{Professio fidei} distinguishes three categories of teaching and corresponding assent. & Governs why every corrected error is not technically heresy; grade A.\\
1998 & Rome; CDF & Writings of Anthony de Mello & A notification identifies positions incompatible with Catholic faith concerning God, Christ, revelation, and religions. & Text-level doctrinal judgment; no inference about every work or subjective culpability; grade A.\\
1999 & Augsburg; Catholic Church and Lutheran World Federation & Sixteenth-century dispute over justification & Joint Declaration and official common statement identify consensus on basic truths of justification and say the historic condemnations do not apply to the partner's teaching as presented there. & Remaining differences and continuing significance of Trent and Lutheran confessions are expressly retained; grade A.\\
2000 & Rome; CDF & Relativist accounts of Christ, revelation, Church, and religions & \emph{Dominus Iesus} reaffirms the definitive and complete revelation in Christ, his unique universal mediation, and the Church's relation to other religions and Christian bodies. & Doctrinal declaration, not a judgment that every non-Catholic person is a formal heretic; grade A.\\
2001 & Rome; CDF & Interpretation of the Rosmini propositions & A doctrinal note explains that the censured propositions do not belong to Rosmini's authentic position when read in the sense clarified by his full works and later doctrine. & A rare formal illustration of afterlife qualification without denying the 1887 act; grade A.\\
2001 & Rome; CDF & Jacques Dupuis's theology of religious pluralism & A notification requires clarifying propositions on Christ, revelation, Church, and religions while expressly not judging subjective intention. & Text-level clarification, not a personal heresy sentence; grade A.\\
2006--2007 & Rome; CDF & Jon Sobrino's Christological works & A notification identifies propositions judged erroneous or dangerous concerning Christ's divinity, Incarnation, self-consciousness, kingdom, and saving death. & Judgment of writings without judgment of subjective intention; grade A.\\
2018 & Rome; CDF and Francis & Analogical ``neo-Gnosticism'' and ``neo-Pelagianism'' & \emph{Placuit Deo} and \emph{Gaudete et exsultate} warn against salvation reduced to inner knowledge or autonomous effort. & The documents expressly use general analogy rather than asserting historical identity or a new juridical sect; grade A.\\
2026 & Rome; DDF prefect & SSPX episcopal ordinations announced without papal mandate & The 13 May statement says the proposed act, if carried out, will constitute a schismatic act and recalls the consequences of formal adherence. & Prospective canonical warning, not proof that the announced ordinations occurred and not a named-heresy judgment; checked through 16 July; grade A.\\
\end{historytimeline}
